%% Example research chapter, structured for a manuscript-based
%% dissertation: this chapter corresponds to one paper, and each
%% section of the paper lives in its own file in this directory.
%%
%% To add a new chapter, copy this directory, adjust the \chapter
%% title and \label, and update the \subimport lines and the
%% corresponding \input in main.tex.

% Chapter-specific commands.  Define macros used ONLY in this chapter
% here; macros shared across chapters belong in macros.tex to avoid
% "already defined" errors.
\newcommand{\toolone}{\textsc{ToolOne}\xspace}

\chapter{Title of Your First Research Chapter}
\label{chap:example-one}

% This chapter is based on:
%   Your Name, Coauthor One, and Coauthor Two. "Title of the First
%   Paper". Venue Name (VENUE 20XX).

% \subimport (from the `import' package) is used instead of \input so
% that relative paths inside the section files -- most importantly
% \includegraphics{figures/...} -- resolve against this chapter's
% directory rather than the repository root.
\subimport{chapter-example-1/}{introduction}
\subimport{chapter-example-1/}{approach}
\subimport{chapter-example-1/}{evaluation}
\subimport{chapter-example-1/}{related-work}
\subimport{chapter-example-1/}{conclusion}
